\section{Cones and termwise split sequences}
\label{section-cones}

\noindent
Let $\mathcal{A}$ be an additive category, and let
$K(\mathcal{A})$ denote the category of complexes of
$\mathcal{A}$ with morphisms given by morphisms of
complexes up to homotopy. Note that the shift functors
$[n]$ on complexes, see
Homology, Definition \ref{homology-definition-shift-cochain},
give rise to functors $[n] : K(\mathcal{A}) \to K(\mathcal{A})$
such that $[n] \circ [m] = [n + m]$ and $[0] = \text{id}$.

\begin{definition}
\label{definition-cone}
Let $\mathcal{A}$ be an additive category.
Let $f : K^\bullet \to L^\bullet$ be a morphism of
complexes of $\mathcal{A}$. The {\it cone} of $f$
is the complex $C(f)^\bullet$ given by
$C(f)^n = L^n \oplus K^{n + 1}$ and
differential
$$
d_{C(f)}^n =
\left(
\begin{matrix}
d^n_L & f^{n + 1} \\
0 & -d_K^{n + 1}
\end{matrix}
\right)
$$
It comes equipped with canonical morphisms of complexes
$i : L^\bullet \to C(f)^\bullet$ and $p : C(f)^\bullet \to K^\bullet[1]$
induced by the obvious maps $L^n \to C(f)^n \to K^{n + 1}$.
\end{definition}

\noindent
In other words $(K, L, C(f), f, i, p)$ forms a triangle:
$$
K^\bullet \to L^\bullet \to C(f)^\bullet \to K^\bullet[1]
$$
The formation of this triangle is
functorial in the following sense.

\begin{lemma}
\label{lemma-functorial-cone}
Suppose that
$$
\xymatrix{
K_1^\bullet \ar[r]_{f_1} \ar[d]_a & L_1^\bullet \ar[d]^b \\
K_2^\bullet \ar[r]^{f_2} & L_2^\bullet
}
$$
is a diagram of morphisms of complexes which is commutative
up to homotopy. Then there exists a morphism
$c : C(f_1)^\bullet \to C(f_2)^\bullet$ which gives rise to
a morphism of triangles
$(a, b, c) : (K_1^\bullet, L_1^\bullet, C(f_1)^\bullet, f_1, i_1, p_1)
\to
(K_2^\bullet, L_2^\bullet, C(f_2)^\bullet, f_2, i_2, p_2)$
of $K(\mathcal{A})$.
\end{lemma}

\begin{proof}
Let $h^n : K_1^n \to L_2^{n - 1}$ be a family of morphisms such that
$b \circ f_1 - f_2 \circ a= d \circ h + h \circ d$.
Define $c^n$ by the matrix
$$
c^n =
\left(
\begin{matrix}
b^n & h^{n + 1} \\
0 & a^{n + 1}
\end{matrix}
\right) :
L_1^n \oplus K_1^{n + 1} \to L_2^n \oplus K_2^{n + 1}
$$
A matrix computation show that $c$ is a morphism of complexes.
It is trivial that $c \circ i_1 = i_2 \circ b$, and it is
trivial also to check that $p_2 \circ c = a \circ p_1$.
\end{proof}

\noindent
Note that the morphism $c : C(f_1)^\bullet \to C(f_2)^\bullet$
constructed in the
proof of Lemma \ref{lemma-functorial-cone} in general depends on the
chosen homotopy $h$ between $f_2 \circ a$ and $b \circ f_1$.

\begin{lemma}
\label{lemma-map-from-cone}
Suppose that $f: K^\bullet \to L^\bullet$ and $g : L^\bullet \to M^\bullet$
are morphisms of complexes such that $g \circ f$ is homotopic to zero.
Then
\begin{enumerate}
\item $g$ factors through a morphism $C(f)^\bullet \to M^\bullet$, and
\item $f$ factors through a morphism $K^\bullet \to C(g)^\bullet[-1]$.
\end{enumerate}
\end{lemma}

\begin{proof}
The assumptions say that the diagram
$$
\xymatrix{
K^\bullet \ar[r]_f \ar[d] & L^\bullet \ar[d]^g \\
0 \ar[r] & M^\bullet
}
$$
commutes up to homotopy.
Since the cone on $0 \to M^\bullet$ is $M^\bullet$ the
map $C(f)^\bullet \to C(0 \to M^\bullet) = M^\bullet$
of Lemma \ref{lemma-functorial-cone}
is the map in (1). The cone on $K^\bullet \to 0$ is
$K^\bullet[1]$ and applying Lemma \ref{lemma-functorial-cone}
gives a map $K^\bullet[1] \to C(g)^\bullet$. Applying
$[-1]$ we obtain the map in (2).
\end{proof}

\noindent
Note that the morphisms $C(f)^\bullet \to M^\bullet$ and
$K^\bullet \to C(g)^\bullet[-1]$ constructed in the proof
of Lemma \ref{lemma-map-from-cone} in general depend on the
chosen homotopy.

\begin{definition}
\label{definition-termwise-split-map}
Let $\mathcal{A}$ be an additive category.
A {\it termwise split injection $\alpha : A^\bullet \to B^\bullet$}
is a morphism of complexes such that each $A^n \to B^n$
is isomorphic to the inclusion of a direct summand.
A {\it termwise split surjection $\beta : B^\bullet \to C^\bullet$}
is a morphism of complexes such that each $B^n \to C^n$
is isomorphic to the projection onto a direct summand.
\end{definition}

\begin{lemma}
\label{lemma-make-commute-map}
Let $\mathcal{A}$ be an additive category.
Let
$$
\xymatrix{
A^\bullet \ar[r]_f \ar[d]_a & B^\bullet \ar[d]^b \\
C^\bullet \ar[r]^g & D^\bullet
}
$$
be a diagram of morphisms of complexes commuting up to homotopy.
If $f$ is a termwise split injection, then $b$ is homotopic to a
morphism which makes the diagram commute.
If $g$ is a termwise split surjection, then $a$ is homotopic to a
morphism which makes the diagram commute.
\end{lemma}

\begin{proof}
Let $h^n : A^n \to D^{n - 1}$ be a collection of morphisms
such that $bf - ga = dh + hd$. Suppose that $\pi^n : B^n \to A^n$
are morphisms splitting the morphisms $f^n$.
Take $b' = b - dh\pi - h\pi d$.
Suppose $s^n : D^n \to C^n$ are morphisms splitting the morphisms
$g^n : C^n \to D^n$. Take $a' = a + dsh + shd$.
Computations omitted.
\end{proof}

\noindent
The following lemma can be used to replace a morphism of complexes
by a morphism where in each degree the map is the injection of a
direct summand.

\begin{lemma}
\label{lemma-make-injective}
Let $\mathcal{A}$ be an additive category.
Let $\alpha : K^\bullet \to L^\bullet$ be a morphism
of complexes of $\mathcal{A}$.
There exists a factorization
$$
\xymatrix{
K^\bullet \ar[r]^{\tilde \alpha} \ar@/_1pc/[rr]_\alpha &
\tilde L^\bullet \ar[r]^\pi &
L^\bullet
}
$$
such that
\begin{enumerate}
\item $\tilde \alpha$ is a termwise split injection (see
Definition \ref{definition-termwise-split-map}),
\item there is a map of complexes $s : L^\bullet \to \tilde L^\bullet$
such that $\pi \circ s = \text{id}_{L^\bullet}$ and such that
$s \circ \pi$ is homotopic to $\text{id}_{\tilde L^\bullet}$.
\end{enumerate}
Moreover, if both $K^\bullet$ and $L^\bullet$ are in
$K^{+}(\mathcal{A})$, $K^{-}(\mathcal{A})$, or $K^b(\mathcal{A})$,
then so is $\tilde L^\bullet$.
\end{lemma}

\begin{proof}
We set
$$
\tilde L^n = L^n \oplus K^n \oplus K^{n + 1}
$$
and we define
$$
d^n_{\tilde L} =
\left(
\begin{matrix}
d^n_L & 0 & 0 \\
0 & d^n_K & \text{id}_{K^{n + 1}} \\
0 & 0 & -d^{n + 1}_K
\end{matrix}
\right)
$$
In other words, $\tilde L^\bullet = L^\bullet \oplus C(1_{K^\bullet})$.
Moreover, we set
$$
\tilde \alpha =
\left(
\begin{matrix}
\alpha \\
\text{id}_{K^n} \\
0
\end{matrix}
\right)
$$
which is clearly a split injection. It is also clear that it defines a morphism
of complexes. We define
$$
\pi =
\left(
\begin{matrix}
\text{id}_{L^n} &
0 &
0
\end{matrix}
\right)
$$
so that clearly $\pi \circ \tilde \alpha = \alpha$. We set
$$
s =
\left(
\begin{matrix}
\text{id}_{L^n} \\
0 \\
0
\end{matrix}
\right)
$$
so that $\pi \circ s = \text{id}_{L^\bullet}$. Finally,
let $h^n : \tilde L^n \to \tilde L^{n - 1}$ be the map
which maps the summand $K^n$ of $\tilde L^n$ via the identity morphism
to the summand $K^n$ of $\tilde L^{n - 1}$. Then it is a trivial matter
(see computations in remark below) to prove that
$$
\text{id}_{\tilde L^\bullet} - s \circ \pi
=
d \circ h + h \circ d
$$
which finishes the proof of the lemma.
\end{proof}

\begin{remark}
\label{remark-compute-modules}
To see the last displayed equality in the proof above we can argue
with elements as follows. We have
$s\pi(l, k, k^{+}) = (l, 0, 0)$.
Hence the morphism of the left hand side maps
$(l, k, k^{+})$ to $(0, k, k^{+})$.
On the other hand $h(l, k, k^{+}) = (0, 0, k)$ and
$d(l, k, k^{+}) = (dl, dk + k^{+}, -dk^{+})$.
Hence $(dh + hd)(l, k, k^{+}) =
d(0, 0, k) + h(dl, dk + k^{+}, -dk^{+}) =
(0, k, -dk) + (0, 0, dk + k^{+}) = (0, k, k^{+})$
as desired.
\end{remark}

\begin{lemma}
\label{lemma-make-surjective}
Let $\mathcal{A}$ be an additive category.
Let $\alpha : K^\bullet \to L^\bullet$ be a morphism
of complexes of $\mathcal{A}$.
There exists a factorization
$$
\xymatrix{
K^\bullet \ar[r]^i \ar@/_1pc/[rr]_\alpha &
\tilde K^\bullet \ar[r]^{\tilde \alpha} &
L^\bullet
}
$$
such that
\begin{enumerate}
\item $\tilde \alpha$ is a termwise split surjection (see
Definition \ref{definition-termwise-split-map}),
\item there is a map of complexes $s : \tilde K^\bullet \to K^\bullet$
such that $s \circ i = \text{id}_{K^\bullet}$ and such that
$i \circ s$ is homotopic to $\text{id}_{\tilde K^\bullet}$.
\end{enumerate}
Moreover, if both $K^\bullet$ and $L^\bullet$ are in
$K^{+}(\mathcal{A})$, $K^{-}(\mathcal{A})$, or $K^b(\mathcal{A})$,
then so is $\tilde K^\bullet$.
\end{lemma}

\begin{proof}
Dual to Lemma \ref{lemma-make-injective}.
Take
$$
\tilde K^n = K^n \oplus L^{n - 1} \oplus L^n
$$
and we define
$$
d^n_{\tilde K} =
\left(
\begin{matrix}
d^n_K & 0 & 0 \\
0 & - d^{n - 1}_L & \text{id}_{L^n} \\
0 & 0 & d^n_L
\end{matrix}
\right)
$$
in other words $\tilde K^\bullet = K^\bullet \oplus C(1_{L^\bullet[-1]})$.
Moreover, we set
$$
\tilde \alpha =
\left(
\begin{matrix}
\alpha &
0 &
\text{id}_{L^n}
\end{matrix}
\right)
$$
which is clearly a split surjection. It is also clear that it defines a
morphism of complexes. We define
$$
i =
\left(
\begin{matrix}
\text{id}_{K^n} \\
0 \\
0
\end{matrix}
\right)
$$
so that clearly $\tilde \alpha \circ i = \alpha$. We set
$$
s =
\left(
\begin{matrix}
\text{id}_{K^n} &
0 &
0
\end{matrix}
\right)
$$
so that $s \circ i = \text{id}_{K^\bullet}$. Finally,
let $h^n : \tilde K^n \to \tilde K^{n - 1}$ be the map
which maps the summand $L^{n - 1}$ of $\tilde K^n$ via the identity morphism
to the summand $L^{n - 1}$ of $\tilde K^{n - 1}$. Then it is a trivial matter
to prove that
$$
\text{id}_{\tilde K^\bullet} - i \circ s
=
d \circ h + h \circ d
$$
which finishes the proof of the lemma.
\end{proof}

\begin{definition}
\label{definition-split-ses}
Let $\mathcal{A}$ be an additive category.
A {\it termwise split exact sequence of complexes of $\mathcal{A}$}
is a complex of complexes
$$
0 \to
A^\bullet \xrightarrow{\alpha}
B^\bullet \xrightarrow{\beta}
C^\bullet \to 0
$$
together with given direct sum decompositions
$B^n = A^n \oplus C^n$
compatible with $\alpha^n$ and $\beta^n$.
We often write $s^n : C^n \to B^n$ and $\pi^n : B^n \to A^n$
for the maps induced by the direct sum decompositions.
According to
Homology, Lemma \ref{homology-lemma-ses-termwise-split-cochain}
we get an associated morphism of complexes
$$
\delta : C^\bullet \longrightarrow A^\bullet[1]
$$
which in degree $n$ is the map $\pi^{n + 1} \circ d_B^n \circ s^n$.
In other words
$(A^\bullet, B^\bullet, C^\bullet, \alpha, \beta, \delta)$
forms a triangle
$$
A^\bullet \to B^\bullet \to C^\bullet \to A^\bullet[1]
$$
This will be the {\it triangle associated to the termwise
split sequence of complexes}.
\end{definition}

\begin{lemma}
\label{lemma-triangle-independent-splittings}
Let $\mathcal{A}$ be an additive category. Let
$0 \to A^\bullet \to B^\bullet \to C^\bullet \to 0$
be termwise split exact sequences as in
Definition \ref{definition-split-ses}.
Let $(\pi')^n$, $(s')^n$ be a second collection of splittings.
Denote $\delta' : C^\bullet \longrightarrow A^\bullet[1]$ the
morphism associated to this second set of splittings.
Then
$$
(1, 1, 1) :
(A^\bullet, B^\bullet, C^\bullet, \alpha, \beta, \delta)
\longrightarrow
(A^\bullet, B^\bullet, C^\bullet, \alpha, \beta, \delta')
$$
is an isomorphism of triangles in $K(\mathcal{A})$.
\end{lemma}

\begin{proof}
The statement simply means that $\delta$ and $\delta'$ are
homotopic maps of complexes. This is
Homology, Lemma \ref{homology-lemma-ses-termwise-split-homotopy-cochain}.
\end{proof}

\begin{remark}
\label{remark-make-commute}
Let $\mathcal{A}$ be an additive category.
Let $0 \to A_i^\bullet \to B_i^\bullet \to C_i^\bullet \to 0$, $i = 1, 2$
be termwise split exact sequences. Suppose that
$a : A_1^\bullet \to A_2^\bullet$,
$b : B_1^\bullet \to B_2^\bullet$, and
$c : C_1^\bullet \to C_2^\bullet$ are morphisms of complexes
such that
$$
\xymatrix{
A_1^\bullet \ar[d]_a \ar[r] &
B_1^\bullet \ar[r] \ar[d]_b &
C_1^\bullet \ar[d]_c \\
A_2^\bullet \ar[r] & B_2^\bullet \ar[r] & C_2^\bullet
}
$$
commutes in $K(\mathcal{A})$. In general, there does {\bf not} exist
a morphism $b' : B_1^\bullet \to B_2^\bullet$ which is homotopic to $b$
such that the diagram above commutes in the category of complexes.
Namely, consider
Examples, Equation (\ref{examples-equation-commutes-up-to-homotopy}).
If we could replace the middle map there by a homotopic one such that
the diagram commutes, then we would have additivity of traces which we do not.
\end{remark}

\begin{lemma}
\label{lemma-nilpotent}
Let $\mathcal{A}$ be an additive category.
Let $0 \to A_i^\bullet \to B_i^\bullet \to C_i^\bullet \to 0$, $i = 1, 2, 3$
be termwise split exact sequences of complexes. Let
$b : B_1^\bullet \to B_2^\bullet$ and $b' : B_2^\bullet \to B_3^\bullet$
be morphisms of complexes such that
$$
\vcenter{
\xymatrix{
A_1^\bullet \ar[d]_0 \ar[r] &
B_1^\bullet \ar[r] \ar[d]_b &
C_1^\bullet \ar[d]_0 \\
A_2^\bullet \ar[r] & B_2^\bullet \ar[r] & C_2^\bullet
}
}
\quad\text{and}\quad
\vcenter{
\xymatrix{
A_2^\bullet \ar[d]^0 \ar[r] &
B_2^\bullet \ar[r] \ar[d]^{b'} &
C_2^\bullet \ar[d]^0 \\
A_3^\bullet \ar[r] & B_3^\bullet \ar[r] & C_3^\bullet
}
}
$$
commute in $K(\mathcal{A})$. Then $b' \circ b = 0$ in $K(\mathcal{A})$.
\end{lemma}

\begin{proof}
By Lemma \ref{lemma-make-commute-map} we can replace $b$ and $b'$ by homotopic
maps such that the right square of the left diagram commutes and the
left square of the right diagram commutes. In other words, we have
$\Im(b^n) \subset \Im(A_2^n \to B_2^n)$ and
$\Ker((b')^n) \supset \Im(A_2^n \to B_2^n)$.
Then $b' \circ b = 0$ as a map of complexes.
\end{proof}

\begin{lemma}
\label{lemma-third-isomorphism}
Let $\mathcal{A}$ be an additive category.
Let $f_1 : K_1^\bullet \to L_1^\bullet$ and
$f_2 : K_2^\bullet \to L_2^\bullet$ be morphisms of complexes.
Let
$$
(a, b, c) :
(K_1^\bullet, L_1^\bullet, C(f_1)^\bullet, f_1, i_1, p_1)
\longrightarrow
(K_2^\bullet, L_2^\bullet, C(f_2)^\bullet, f_2, i_2, p_2)
$$
be any morphism of triangles of $K(\mathcal{A})$.
If $a$ and $b$ are homotopy equivalences then so is $c$.
\end{lemma}

\begin{proof}
Let $a^{-1} : K_2^\bullet \to K_1^\bullet$ be a morphism of complexes which
is inverse to $a$ in $K(\mathcal{A})$.
Let $b^{-1} : L_2^\bullet \to L_1^\bullet$ be a morphism of complexes which
is inverse to $b$ in $K(\mathcal{A})$.
Let $c' : C(f_2)^\bullet \to C(f_1)^\bullet$
be the morphism from Lemma \ref{lemma-functorial-cone} applied
to $f_1 \circ a^{-1} = b^{-1} \circ f_2$. If we can show that
$c \circ c'$ and $c' \circ c$ are isomorphisms in $K(\mathcal{A})$
then we win. Hence it suffices to prove the following: Given
a morphism of triangles
$(1, 1, c) : (K^\bullet, L^\bullet, C(f)^\bullet, f, i, p)$
in $K(\mathcal{A})$ the morphism $c$ is an isomorphism in $K(\mathcal{A})$.
By assumption the two squares in the diagram
$$
\xymatrix{
L^\bullet \ar[r] \ar[d]_1 &
C(f)^\bullet \ar[r] \ar[d]_c &
K^\bullet[1] \ar[d]_1 \\
L^\bullet \ar[r] &
C(f)^\bullet \ar[r] &
K^\bullet[1]
}
$$
commute up to homotopy. By construction of $C(f)^\bullet$ the rows
form termwise split sequences of complexes. Thus we see that
$(c - 1)^2 = 0$ in $K(\mathcal{A})$ by Lemma \ref{lemma-nilpotent}.
Hence $c$ is an isomorphism in $K(\mathcal{A})$ with inverse $2 - c$.
\end{proof}

\noindent
Hence if $a$ and $b$ are homotopy equivalences then
the resulting morphism of triangles is an isomorphism of triangles
in $K(\mathcal{A})$.
It turns out that the collection of triangles of $K(\mathcal{A})$
given by cones and the collection of triangles of $K(\mathcal{A})$
given by termwise split sequences of complexes are the same
up to isomorphisms, at least up to sign!

\begin{lemma}
\label{lemma-the-same-up-to-isomorphisms}
Let $\mathcal{A}$ be an additive category.
\begin{enumerate}
\item Given a termwise split sequence of complexes
$(\alpha : A^\bullet \to B^\bullet,
\beta : B^\bullet \to C^\bullet, s^n, \pi^n)$
there exists a homotopy equivalence $C(\alpha)^\bullet \to C^\bullet$
such that the diagram
$$
\xymatrix{
A^\bullet \ar[r] \ar[d] & B^\bullet \ar[d] \ar[r] &
C(\alpha)^\bullet \ar[r]_{-p} \ar[d] & A^\bullet[1] \ar[d] \\
A^\bullet \ar[r] & B^\bullet \ar[r] &
C^\bullet \ar[r]^\delta & A^\bullet[1]
}
$$
defines an isomorphism of triangles in $K(\mathcal{A})$.
\item Given a morphism of complexes $f : K^\bullet \to L^\bullet$
there exists an isomorphism of triangles
$$
\xymatrix{
K^\bullet \ar[r] \ar[d] & \tilde L^\bullet \ar[d] \ar[r] &
M^\bullet \ar[r]_{\delta} \ar[d] & K^\bullet[1] \ar[d] \\
K^\bullet \ar[r] & L^\bullet \ar[r] &
C(f)^\bullet \ar[r]^{-p} & K^\bullet[1]
}
$$
where the upper triangle is the triangle associated to a
termwise split exact sequence $K^\bullet \to \tilde L^\bullet \to M^\bullet$.
\end{enumerate}
\end{lemma}

\begin{proof}
Proof of (1). We have $C(\alpha)^n = B^n \oplus A^{n + 1}$
and we simply define $C(\alpha)^n \to C^n$ via the projection
onto $B^n$ followed by $\beta^n$. This defines
a morphism of complexes because the compositions
$A^{n + 1} \to B^{n + 1} \to C^{n + 1}$ are zero.
To get a homotopy inverse we take
$C^\bullet \to C(\alpha)^\bullet$ given by
$(s^n , -\delta^n)$ in degree $n$. This is a morphism of complexes
because the morphism $\delta^n$ can be characterized as the
unique morphism $C^n \to A^{n + 1}$ such that
$d \circ s^n - s^{n + 1} \circ d = \alpha \circ \delta^n$,
see proof of
Homology, Lemma \ref{homology-lemma-ses-termwise-split-cochain}.
The composition
$C^\bullet \to C(\alpha)^\bullet \to C^\bullet$ is the identity.
The composition $C(\alpha)^\bullet \to C^\bullet \to C(\alpha)^\bullet$
is equal to the morphism
$$
\left(
\begin{matrix}
s^n \circ \beta^n & 0 \\
-\delta^n \circ \beta^n & 0
\end{matrix}
\right)
$$
To see that this is homotopic to the identity map
use the homotopy $h^n : C(\alpha)^n \to C(\alpha)^{n - 1}$
given by the matrix
$$
\left(
\begin{matrix}
0 & 0 \\
\pi^n & 0
\end{matrix}
\right) : C(\alpha)^n = B^n \oplus A^{n + 1} \to
B^{n - 1} \oplus A^n = C(\alpha)^{n - 1}
$$
It is trivial to verify that
$$
\left(
\begin{matrix}
1 & 0 \\
0 & 1
\end{matrix}
\right)
-
\left(
\begin{matrix}
s^n \\
-\delta^n
\end{matrix}
\right)
\left(
\begin{matrix}
\beta^n & 0
\end{matrix}
\right)
=
\left(
\begin{matrix}
d & \alpha^n \\
0 & -d
\end{matrix}
\right)
\left(
\begin{matrix}
0 & 0 \\
\pi^n & 0
\end{matrix}
\right)
+
\left(
\begin{matrix}
0 & 0 \\
\pi^{n + 1} & 0
\end{matrix}
\right)
\left(
\begin{matrix}
d & \alpha^{n + 1} \\
0 & -d
\end{matrix}
\right)
$$
To finish the proof of (1) we have to show that the morphisms
$-p : C(\alpha)^\bullet \to A^\bullet[1]$ (see
Definition \ref{definition-cone})
and  $C(\alpha)^\bullet \to C^\bullet \to A^\bullet[1]$ agree up
to homotopy. This is clear from the above. Namely, we can use the homotopy
inverse $(s, -\delta) : C^\bullet \to C(\alpha)^\bullet$
and check instead that the two maps
$C^\bullet \to A^\bullet[1]$ agree. And note that
$p \circ (s, -\delta) = -\delta$ as desired.

\medskip\noindent
Proof of (2). We let $\tilde f : K^\bullet \to \tilde L^\bullet$,
$s : L^\bullet \to \tilde L^\bullet$
and $\pi : \tilde L^\bullet \to L^\bullet$ be as in
Lemma \ref{lemma-make-injective}. By
Lemmas \ref{lemma-functorial-cone} and \ref{lemma-third-isomorphism}
the triangles $(K^\bullet, L^\bullet, C(f), i, p)$ and
$(K^\bullet, \tilde L^\bullet, C(\tilde f), \tilde i, \tilde p)$
are isomorphic. Note that we can compose isomorphisms of
triangles. Thus we may replace $L^\bullet$ by
$\tilde L^\bullet$ and $f$ by $\tilde f$. In other words
we may assume that $f$ is a termwise split injection.
In this case the result follows from part (1).
\end{proof}


\begin{lemma}
\label{lemma-sequence-maps-split}
Let $\mathcal{A}$ be an additive category.
Let $A_1^\bullet \to A_2^\bullet \to \ldots \to A_n^\bullet$
be a sequence of composable morphisms of complexes.
There exists a commutative diagram
$$
\xymatrix{
A_1^\bullet \ar[r] &
A_2^\bullet \ar[r] &
\ldots \ar[r] &
A_n^\bullet \\
B_1^\bullet \ar[r] \ar[u] &
B_2^\bullet \ar[r] \ar[u] &
\ldots \ar[r] &
B_n^\bullet \ar[u]
}
$$
such that each morphism $B_i^\bullet \to B_{i + 1}^\bullet$
is a termwise split injection and each $B_i^\bullet \to A_i^\bullet$
is a homotopy equivalence. Moreover, if all $A_i^\bullet$ are in
$K^{+}(\mathcal{A})$, $K^{-}(\mathcal{A})$, or $K^b(\mathcal{A})$,
then so are the $B_i^\bullet$.
\end{lemma}

\begin{proof}
The case $n = 1$ is without content.
Lemma \ref{lemma-make-injective} is the case $n = 2$.
Suppose we have constructed the diagram
except for $B_n^\bullet$. Apply Lemma \ref{lemma-make-injective} to
the composition $B_{n - 1}^\bullet \to A_{n - 1}^\bullet \to A_n^\bullet$.
The result is a factorization
$B_{n - 1}^\bullet \to B_n^\bullet \to A_n^\bullet$
as desired.
\end{proof}

\begin{lemma}
\label{lemma-rotate-triangle}
Let $\mathcal{A}$ be an additive category. Let
$(\alpha : A^\bullet \to B^\bullet, \beta : B^\bullet \to C^\bullet, s^n,
\pi^n)$ be a termwise split sequence of complexes.
Let $(A^\bullet, B^\bullet, C^\bullet, \alpha, \beta, \delta)$
be the associated triangle.
Then the triangle
$(C^\bullet[-1], A^\bullet, B^\bullet, \delta[-1], \alpha, \beta)$
is isomorphic to the triangle
$(C^\bullet[-1], A^\bullet, C(\delta[-1])^\bullet, \delta[-1], i, p)$.
\end{lemma}

\begin{proof}
We write $B^n = A^n \oplus C^n$ and we identify $\alpha^n$ and $\beta^n$
with the natural inclusion and projection maps. By construction of $\delta$ we
have
$$
d_B^n =
\left(
\begin{matrix}
d_A^n & \delta^n \\
0 & d_C^n
\end{matrix}
\right)
$$
On the other hand the cone of $\delta[-1] : C^\bullet[-1] \to A^\bullet$
is given as $C(\delta[-1])^n = A^n \oplus C^n$ with differential identical
with the matrix above! Whence the lemma.
\end{proof}

\begin{lemma}
\label{lemma-rotate-cone}
Let $\mathcal{A}$ be an additive category.
Let $f : K^\bullet \to L^\bullet$ be a morphism of complexes.
The triangle $(L^\bullet, C(f)^\bullet, K^\bullet[1], i, p, f[1])$ is
the triangle associated to the termwise split sequence
$$
0 \to L^\bullet \to C(f)^\bullet \to K^\bullet[1] \to 0
$$
coming from the definition of the cone of $f$.
\end{lemma}

\begin{proof}
Immediate from the definitions.
\end{proof}






